% ============================================================
\chapter{事前分布：均衡リターン、リスク回避度、共分散}\label{ch:prior}

BL の出発点は、ビューがないときに持つべき比率と、そのときの期待リターンである。この章では、均衡リターン $\Pi$ の作り方、リスク回避度 $\delta$ の決め方、共分散 $\Sigma$ の推定を扱う。ビューの話に入る前に、ここを固めておく理由は単純で、$\Pi$ と $\Sigma$ がずれていると、どんなにうまくビューを作っても、出発点からずれた比率になるからである。

\section{逆最適化}

平均分散最適化~\eqref{eq:mvobj} の制約のない答えは $w^*=(\delta\Sigma)^{-1}\mu$ である。これを逆に解いて、「$\wmkt$ が答えになるような $\mu$」を求めたのが均衡リターン $\Pi=\delta\Sigma\wmkt$ である（Sharpe 1974 の「比率から期待リターンを推し量る」考え方）。

均衡リターンには、次の性質がある。

\begin{enumerate}
\item $\Pi$ を最適化に入れると、比率は $\wmkt$ に戻る。同じ $\Sigma$ と $\delta$ を使う限り、$\Sigma$ の推定に誤りがあっても、ビューのないときの比率は $\wmkt$ のままである。$\Sigma$ の誤りは、$\Pi$ と最適化の両方に同じように入り、打ち消し合う。
\item 市場全体の期待リターンは $\wmkt\T\Pi=\delta\,\wmkt\T\Sigma\wmkt=\delta\sigma_m^2$ である。$\delta$ を決めることは、市場のリスクプレミアムを決めることと同じである。
\item 資産ごとには $\Pi=\beta\cdot(\delta\sigma_m^2)$ と書ける。$\beta=\Sigma\wmkt/\sigma_m^2$ は、市場ポートフォリオに対する各資産のベータである。つまり $\Pi$ は CAPM の期待リターンである。
\end{enumerate}

1 番目の性質は重要である。BL の比率が市場からずれるのは、ビューがあるときだけである。逆に言うと、ビューがないのに比率が市場からずれていたら、$\Pi$ と最適化で使う $\Sigma$ や $\delta$ が食い違っているか、制約が効いている。本プロジェクトの元の実装では、均衡リターンの単位が誤っていた（日次のリターンから年率の金利を引いていた）が、ビューが 99.9\% 採用されていたため、結果にはほとんど影響しなかった（『総合報告書』問題 5）。均衡がどうであっても結果が変わらないこと自体が、BL が働いていない証拠だった。

\section{リスク回避度 $\delta$ の決め方}

$\delta$ は、市場ポートフォリオの期待超過リターンを、その分散で割ったものである。
\begin{equation}
\delta=\frac{\E[r_m]-r_f}{\sigma_m^2}
\label{eq:delta}
\end{equation}
Idzorek（2007）の例では、市場のリスクプレミアムを 3\%、市場の分散を $\wmkt\T\Sigma\wmkt=\V{idz.mktvar}$（標準偏差 \V{idz.mktsd}\%）として、$\delta\approx3.07$ を得ている。この資料で論文の表から逆算すると $\delta=\V{idz.lam}$ で、均衡リターンは論文の表と最大 \V{idz.pidiff}\% の差で一致した。He と Litterman（1999）は $\delta=2.5$ を使い、本プロジェクトも $\delta=2.5$ を使っている。PyPortfolioOpt にも、市場の価格から式~\eqref{eq:delta} で $\delta$ を求める関数がある。

\begin{tipbox}[コツ：$\delta$ は期間によらない]
式~\eqref{eq:delta} の分子と分母は、どちらも期間に比例して大きくなる（年率なら日次の約 252 倍）。だから、日次のリターンで計算しても年率で計算しても、$\delta$ は同じになる。$\delta=2.5$ は、日次の $\Sigma$ に使っても年率の $\Sigma$ に使ってもよい。ずれるのは、$\Pi$ と $\Sigma$ と $Q$ と $\Omega$ の単位がそろっていないときである。
\end{tipbox}

$\delta$ を大きくすると、$\Pi$ 全体が比例して大きくなる。$Q$ を固定したまま $\delta$ を変えると、驚き $Q-P\Pi$ が変わるので、BL の結果も変わる。$\delta$ を変えたら、$Q$ の決め方（とくに $\Pi$ からの差として決めている場合）も見直す必要がある。

\section{回帰のベータと、黙示のベータ}

BL の文献では、$\Pi$ を「CAPM のリターン」と呼ぶことが多い。Idzorek は、これが混乱のもとになると注意している。ふつうの CAPM のリターンは、過去のリターンを回帰して得たベータを使う。$\Pi$ は、時価総額の比率と共分散から計算した「黙示のベータ」を使う。両者は違う。

Idzorek の例では、世界の市場の指数（UBS GSMI）に対する回帰のベータで作った期待リターンと、$\Pi$ の相関は 99.8\% だった。それでも、それぞれを最適化に入れた比率の相関は 66\% にとどまり、外国債券の比率は大きく違った。期待リターンのわずかな違いが比率を大きく動かす、平均分散最適化の性質がそのまま出ている。BL で $\Pi$ を使うのは、比率を $\wmkt$ に戻すためである。回帰のベータの CAPM を使うと、この性質が失われる。

\section{何を「市場」とするか}

$\wmkt$ に何を使うかも決める必要がある。

\begin{itemize}
\item \textbf{投資対象の中での時価総額の比率}：本プロジェクトは、判断日の \SP{} の構成銘柄のうち、ビューと株価がそろう銘柄の時価総額を合計 1 にそろえて使っている。
\item \textbf{浮動株調整後の比率}：指数の実際の比率に近い。Idzorek は代わりの候補として挙げている。
\item \textbf{ベンチマークの比率}：運用の目標がベンチマークなら、その比率を使うのが自然である。BL の結果は、ベンチマークからの能動的な傾きとして読める（\chapref{ch:te}）。
\item \textbf{等金額の比率}：等金額を出発点にすれば、$\Pi$ は「等金額を最適にする期待リターン」になる。
\item \textbf{ESG などで調整した比率}：Alpern と Rachev（2025）は、ダウ平均の比率と ESG の点数に比例する比率を $w_{\mathrm{eq}}=(1-\lambda)w_C+\lambda w_\xi$ で混ぜている（\chapref{ch:esg}）。
\end{itemize}

どれを使っても計算はできるが、意味が変わる。$\wmkt$ を変えることは、「ビューがないときに持つ比率」を変えることである。

\section{均衡は「縮小推定の目標」でもある}

Meucci（2010）は、BL の事前分布を縮小推定の一種として説明している。過去の平均リターン $\hat\mu$ は非常に不正確な推定値なので、何かの目標 $\pi_0$ に向かって縮める。
\begin{equation}
\mu(s)=(1-s)\hat\mu+s\,\pi_0,\qquad 0\le s\le 1
\end{equation}
BL の事前分布は、目標を均衡 $\Pi$ にして、$s=1$（過去の平均をまったく使わない）にした極端な場合とみなせる。過去の平均の標準誤差は $\Sigma/T$ なので、事前分布 $\mathcal N(\Pi,\tau\Sigma)$ と比べると $\tau\approx1/T$ が 1 つの目安になる（\chapref{ch:tau}）。

この見方をとると、均衡以外の目標（ゼロ、全資産の平均、ファクターモデルの予想）も自然に考えられる。Black と Litterman 自身も、早い段階で目標を変えた版を検討していた。アクティブ運用では、事前の期待超過リターンを 0（ベンチマークに対して中立）と置くのが自然である（\chapref{ch:te}）。

\section{共分散 $\Sigma$ の推定}\label{sec:cov}

$\Sigma$ は BL の 3 か所（均衡、混ぜ方、最適化）で使う。推定の方法はいくつかある。

\begin{description}[leftmargin=2em,style=nextline]
\item[標本共分散] 過去 $T$ 日のリターンの共分散。資産の数 $N$ がデータの数 $T$ 以上だと、逆行列が存在しない。\SP{} の約 500 銘柄に対して 1 年（約 252 日）は足りない。元の実装は、約 20 日分で約 490 銘柄の共分散を作っていた（『総合報告書』問題 6）。
\item[縮小推定] 標本共分散 $S$ を単純な構造の行列 $F$ に寄せる：$\hat\Sigma=(1-a)S+aF$。Ledoit と Wolf（2004）の方法は、寄せる強さ $a$ をデータから決める。本プロジェクトの「直した計算」はこれを使っている。
\item[ファクターモデル] $r=Bf+\epsilon$ と置くと $\Sigma=BFB\T+D$（$F$ はファクターの共分散、$D$ は個別の分散の対角行列）。推定するものが減り、逆行列も安定する。Kolm と Ritter、Spears ほか（2023）の BL--APT はこの形である。この資料のセクター ETF の例では、Fama--French の 6 ファクターが分散の平均 \V{fac.share}\% を説明した（\chapref{ch:esg}）。
\item[指数加重] 新しいデータに大きな重みを付ける。Meucci の説明では、BL の $\Sigma$ は過去のリターンの指数平滑で推定されている。
\item[時変の共分散] GARCH などで、明日の共分散を予想する。Sun・Wu・De（2023）は年金の資産配分で、VARMA と GARCH から長期の共分散を作り、従来の BL よりシャープレシオが改善したと報告している（\chapref{ch:macro}）。
\end{description}

\begin{tipbox}[コツ：$\Sigma$ は「どこで使うか」で選ぶ]
\begin{itemize}
\item 均衡と最適化で同じ $\Sigma$ を使えば、ビューのないときの比率は必ず $\wmkt$ になる。この一貫性は崩さない。
\item 混ぜ方（$\tau\Sigma$ と $\Omega$ の比）でも同じ $\Sigma$ を使う。$\Omega$ を $\Sigma$ に比例させる場合（\chapref{ch:hlmeucci}）、ここがずれると採用率がずれる。
\item ボラティリティが急に変わる時期には、$\Sigma$ を長い期間で推定し、$\Omega$ を短い期間のボラティリティで伸び縮みさせると、採用率が大きく動く（\chapref{ch:macro}）。意図してそうするのでなければ、両方を同じ速さで動かす。
\end{itemize}
\end{tipbox}

\section{単位をそろえる}

本プロジェクトのように、ビューが「翌月の日次の平均リターン」のとき、すべてを日次でそろえる。

\begin{center}
\small
\begin{tabular}{p{3.2cm}p{10.4cm}}
\toprule
量 & 単位 \\
\midrule
$\Sigma$ & 日次のリターンの共分散（過去 252 取引日） \\
$\Pi$ & $\delta\Sigma\wmkt$。日次の超過リターン \\
$Q$ & LLM の予想（日次の平均リターン）から、日次に直した無リスク金利を引いたもの \\
$\Omega$ & 日次の平均リターンの不確かさの分散 \\
$\tau$ & 単位なし（$\tau\Sigma$ が日次の平均リターンの不確かさの分散） \\
\bottomrule
\end{tabular}
\end{center}

ここで 1 つ問題がある。$\Sigma$ は「1 日のリターンの散らばり」だが、$Q$ は「1 か月の日次リターンの平均」の予想である。1 か月（約 21 日）の平均の散らばりは、1 日の散らばりの $1/21$ である。$\tau\Sigma$ をどう解釈するかで、$\Omega$ の適切な大きさが変わる。この点は、ビューが期待値についてのものか実現値についてのものかという問題と重なるので、\chapref{ch:hlmeucci} と \chapref{ch:accuracy} で詳しく扱う。

\begin{keybox}[この章のまとめ]
\begin{itemize}
\item 均衡リターン $\Pi=\delta\Sigma\wmkt$ は、黙示のベータによる CAPM の期待リターンである。同じ $\Sigma$ で最適化すれば、ビューのないときの比率は $\wmkt$ に戻る。
\item $\delta$ は市場のリスクプレミアム ÷ 市場の分散で、期間によらない。2.5〜3 程度がよく使われる。
\item 均衡は、過去の平均を縮める目標と見ることもできる。この見方から $\tau\approx1/T$ という目安が出てくる。
\item $\Sigma$ は 3 か所で同じものを使い、単位（日次・月次・年率）をそろえる。
\end{itemize}
\end{keybox}
